\subsubsection{\bsq{switch} and Pattern Matching}\label{sec:PatternMatchingSwitch}
Pattern matching\index{pattern matching} (refer to \autocite{ORACLE_DOC_PATTERNMATCHING}\index{pattern matching}) is covered in more detail in chapter \tqfullvref{sec:PatternMatching}. Here I want to discuss only the basic formatting of the \lstinline|switch|\index{switch} when used with pattern matching.

\keeplisting{14}
A \lstinline|switch| statement\index{switch!statement} for the ``\verb#case L:#'' variant would look like this:
\index{switch!case}\index{switch!default}\index{switch!break}
\begin{lstlisting}
var selector = <anObject>
switch( selector )
{
    case <class1> <name1>: <singleStatement>;

    case <class2> <name2>:
    {
        <statements>;
    }

    default: <singleStatement>;
}
\end{lstlisting}

\keeplisting{15}
For the ``\verb#case L:#'' variant, a \lstinline|switch| statement\index{switch!statement} is written as below:
\index{switch!case}\index{switch!default}
\begin{lstlisting}
var selector = <anObject>
switch( selector )
{
    case null -> <singleStatement>;
    
    case <class1> <name1> -> <singleStatement>;

    case <class2> <name2> ->
    {
        <statements>;
    }

    default -> <singleStatement>;
}
\end{lstlisting}

\keeplisting{16}
And an ``\verb#case L->#'' variant \lstinline|switch| expression\index{switch!expression} would have this form:
\index{switch!case}\index{switch!default}\index{switch!yield}
\begin{lstlisting}
var selector = <anObject>
var result = switch( selector )
{
    case null -> <expression>;

    case <class1> <name1> -> <expression>;

    case <class2> <name2> ->
    {
        <expressions>;
        yield <value>;
    }

    default -> <expression>;
}
\end{lstlisting}

Same as for the \lstinline|default|\index{switch!default} branch, the \lstinline|case null| \index{switch!case null} branch can be a statement block or a complex expression.

And an expression can also be always a \lstinline|throw|\index{throw} statement.



%==============================================================================
% End of file!!
%==============================================================================
